\subsection{双代数中的除幂}

令 A 为有单位 k-代数，\(x \in A\)，\(d \in k\)，\(n \in N\)。置

\[
x ^ {(n, d)} = \frac {x (x - d) \dots (x - d (n - 1))}{n !} = \prod_ {i = 0} ^ {n - 1} (x - i d) / (i + 1).\tag{1}
\]

特别地，\(x^{(0,d)} = 1\)，\(x^{(1,d)} = x\)。约定当 n 为小于 0 的整数时 \(x^{(n,d)} = 0\)。置\(n\)\(< 0\)

\[
x ^ {(n)} = x ^ {(n, 0)} = \frac {x ^ {n}}{n !}\tag{2}
\]

\[
\binom {x} {n} = x ^ {(n, 1)} = \frac {x (x - 1) \ldots (x - n + 1)}{n !}.\tag{3}
\]

命题 1。令 A 为双代数，余积为 \(c\)，\(x\) 为 A 的一个本原元（第 II 章第1节第2号）。则

\[
c (x ^ {(n, d)}) = \sum_ {p \in \mathbf {N}, q \in \mathbf {N}, p + q = n} x ^ {(p, d)} \otimes x ^ {(q, d)}.\tag{4}
\]

当 \(n \leq 0\) 时命题显然成立。我们对 \(n\) 作归纳。若公式 (4) 对 \(n\) 成立，则

\[
\begin{array}{l} (n + 1) c (x ^ {(n + 1, d)}) = c (x - d n) c (x ^ {(n, d)}) \\ \quad = (x \otimes 1 + 1 \otimes x - d n   1 \otimes 1) c (x ^ {(n, d)}) \\ \quad = \sum_ {p + q = n} [ x x ^ {(p, d)} \otimes x ^ {(q, d)} + x ^ {(p, d)} \otimes x x ^ {(q, d)} - (p + q) d x ^ {(p, d)} \otimes x ^ {(q, d)} ] \\ \quad = \sum_ {p + q = n} (x - p d) x ^ {(p, d)} \otimes x ^ {(q, d)} + \sum_ {p + q = n} x ^ {(p, d)} \otimes (x - q d) x ^ {(q, d)} \\ \quad = \sum_ {p + q = n} (p + 1) x ^ {(p + 1, d)} \otimes x ^ {(q, d)} + \sum_ {p + q = n} (q + 1) x ^ {(p, d)} \otimes x ^ {(q + 1, d)} \\ \quad = \sum_ {r + s = n + 1} r x ^ {(r, d)} \otimes x ^ {(s, d)} + \sum_ {r + s = n + 1} s x ^ {(r, d)} \otimes x ^ {(s, d)} \\ \quad = (n + 1) \sum_ {r + s = n + 1} x ^ {(r, d)} \otimes x ^ {(s, d)}, \end{array}
\]

因而公式 (4) 对 \(n + 1\) 成立。

